\documentclass{article}
\usepackage{graphicx} % Required for inserting images
\usepackage{amsmath, amssymb}
\usepackage[0T1]


\title{Apêndice - Lógica de primeira ordem}
\author{106618 }
\date{}
\begin{document}

\maketitle

\textbf{Sintaxe}: 
\begin{itemize}
    \item \underline{Assinatura}: $\Sigma = (\mathcal{F}, \mathcal{P}, \tau)$, onde $\mathcal{F} \;\text{e}\;\mathcal{P}$ representam, o conjunto de símbolos de funções e conjunto de símbolos de predicados/relações, respetivamente. Além disso, $\tau: \mathcal{F} \cup \mathcal{P} \rightarrow \mathbb{N}$ representa a aridade dos símbolos anteriores e definimos $A_n= \tau^{-1}(n) \cap A$, para $A=\mathcal{F,P}$
    
    \item \underline{Variáveis}: Conjunto $X$ de símbolos de variáveis.\\

    \item  \underline{Termos}: Conjunto $T_\Sigma$ de símbolos de termos definidos indutivamente por:
       
        \begin{itemize}
            \renewcommand{\labelitemi}{-}
            \item $x\in X$, então $x \in T_\Sigma$
            \item $c \in \mathcal{F}_0$, então $c \in T_\Sigma$
            \item $f \in \mathcal{F}_n$ e $t_1,...,t_n \in T_\Sigma$, então $f(t_1,...,t_n) \in T_\Sigma$
        \end{itemize}

    \item \underline{Fórmulas} Conjunto $L_\Sigma$ de fórmulas definido indutivamente por:

        \begin{itemize}
            \renewcommand{\labelitemi}{-}
            \item $p \in \mathcal{P}_n$ e $t_1,...,t_n \in T_\Sigma$, então $p(t_1,...,t_n) \in L_\Sigma$
            \item $A,B \in L_\Sigma$, então $(\neg A) \in L_\Sigma$ e $(A \implies B) \in L_\Sigma$
            \item $A \in L_\Sigma, x\in X$, então $(\forall xA) \in L_\Sigma$
        \end{itemize}
\end{itemize}

\vspace{1cm}

\textbf{Semântica}
\begin{itemize}
    \item \underline{Estrutura de Interpretação}: $I=(\mathcal{D}, \_{^\mathcal{F}},\_^\mathcal{P})$, onde $\mathcal{D}$ é um conjunto não vazio designado domínio. Além disso, $-^\mathcal{F} \text{ e } -^\mathcal{P}$ são as interpretações de função $\mathcal{F}$ e interpretações de predicado $\mathcal{P}$, respetivamente, definidas como:

    \item \underline{Atribuição de valor ás variáveis}: $\rho: X \rightarrow\mathcal{D}$\\
    Além disso, definimos $\rho(x:= d)(y)=
        \begin{cases}
            \rho(y),\;y\neq x\\
            d\;\;\;\,\;,\;y=x
        \end{cases}$

    \item \underline{Interpretação de termos}: $[|\_|]_{I_\rho}: T_\Sigma \rightarrow \mathcal{D}$
        \begin{itemize}
            \renewcommand{\labelitemi}{-}
            \item $[|x|]_{I_\rho}=\rho(x)$
            \item $[|c|]_{I_\rho}=\underline{c}^\mathcal{F}$
            \item $[|f(t_1,...,t_n)|]_{I_\rho}=\underline{f}^\mathcal{F}([|t_1|]_{I_\rho},...,[|t_n|]_{I_\rho})$
        \end{itemize}

    \item \underline{Interpretação de fórmulas}: $[|\_|]_{I_\rho}: L_\Sigma \rightarrow\{0,1\}$

    \item \underline{Satisfação local de fórmulas}: Fixados $I$ e $\rho$, 
        \begin{itemize}
            \renewcommand{\labelitemi}{-}
            \item $I_\rho \Vdash p(t_1,...,t_n)$ sse $([|t_1|]_{I_\rho},...,[|t_n|]_{I_\rho}) \in \underline{p}^\mathcal{P}$
            \item $I_\rho \Vdash (\neg A)$ sse $I_\rho \nVdash A$ 
            \item $I_\rho \Vdash (A \implies B)$ sse $I_\rho \nVdash A$ ou $I_\rho \Vdash B$
            \item $I_\rho \Vdash  (\forall x A)$ sse $I_{\rho(x:=d)} \Vdash A$ para qualquer $d \in \mathcal{D}$
        \end{itemize}

    \item \underline{Satisfação global de fórmulas}:\\
    $I \Vdash A$ sse $I_\rho \Vdash A$ qualquer que seja a atribuição $\rho$.

    \item \underline{Consequência semântica}: Relação $\models \; \subseteq \mathcal{P}(L_\Sigma) \times L_\Sigma$ definida como:\\
    $\Gamma \models A $ sse para qualquer estrutura $I$ tal que se $I \Vdash B$ para todo o $B \in \Gamma$, então $I \Vdash A$
\end{itemize}


\end{document}
